\section{Umfang des MVP}
\label{chap:mvp}

\subsection{Kleinste nützliche Version}

Die kleinste nützliche Version der App ermöglicht es, eine Hautstelle zu fotografieren, eine Risikoeinschätzung zu erhalten und daraus eine klare Handlungsempfehlung abzuleiten. Schon damit erfüllt die App ihren Kernzweck: Nutzer erkennen, ob eine Hautveränderung ärztlich abgeklärt werden sollte.

\subsection{MVP-Funktionalitäten}
\label{sec:mvp-funktionen}

Aus der Marktanalyse~\cite{marktanalyse2026} haben wir die folgenden Funktionalitäten für das Minimum Viable Product (MVP) ausgewählt. Die zugehörigen Anwendungsfälle sind in \chapref{usecases} beschrieben.

\begin{table}[htbp]
\caption{MVP-Funktionalitäten}
\label{tab:mvp-features}
\footnotesize
\renewcommand{\arraystretch}{1.15}
\begin{tabularx}{\textwidth}{@{}l l >{\RaggedRight\arraybackslash}p{2.9cm} >{\RaggedRight\arraybackslash}X >{\RaggedRight\arraybackslash}p{2.2cm} >{\RaggedRight\arraybackslash}p{1.7cm}@{}}
\toprule
\textbf{ID} & \textbf{Bereich} & \textbf{Feature} & \textbf{Beschreibung} & \textbf{Gesehen bei} & \textbf{Use Case} \\
\midrule
F1 & Bildaufnahme & Geführte Kameraaufnahme & Overlay für Abstand und Bildausschnitt, manueller Auslöser & SkinVision~\cite{app-skinvision}, SkinScreener~\cite{app-skinscreener}, DERM Zero~\cite{app-dermzero} & UC-03 \\
F2 & Bildaufnahme & Bildqualitätsprüfung & Schärfe und Belichtung nach der Aufnahme prüfen und bei Mängeln eine Neuaufnahme verlangen & SkinVision, DermEngine~\cite{app-dermengine} & UC-04 \\
F3 & KI-Analyse & Risikoeinstufung & 3 Stufen (niedrig / mittel / hoch, Ampel) & SkinVision, SkinScreener, Skinive~\cite{app-skinive} & UC-05 \\
F4 & KI-Analyse & Ausschluss ungeeigneter Fälle & Kinder, Hauttyp V--VI, Schleimhaut, Nägel $\rightarrow$ direkt zum Arzt & SkinScreener (Hauttyp I--IV), NICE-Auflage bei DERM~\cite{app-derm} & UC-02, UC-07 \\
F5 & Ergebnis & Handlungsempfehlung & \enquote{In 3 Monaten erneut prüfen} vs. \enquote{Dringend einen Arzt aufsuchen}, immer mit dem Hinweis, dass das Ergebnis nur eine Einschätzung ist, keinen ärztlichen Rat ersetzt und nur ein Arzt eine Diagnose stellen kann & SkinVision, SkinScreener, DERM Zero & UC-06 \\
F6 & Monitoring & Läsionsprofil und Verlauf & Jede Hautstelle mit eigener Bildhistorie, Aufnahmen und Hautstellen lassen sich löschen & SkinVision, Miiskin~\cite{app-miiskin}, MoleScope~\cite{app-molescope} & UC-08, UC-09, UC-12 \\
F7 & Monitoring & 2D-Bodymap & Läsionen auf einer Körperübersicht markieren und anzeigen & SkinVision, MoleScope, DermEngine & UC-08, UC-09 \\
F8 & Engagement & Erinnerungen & Lokale Benachrichtigung auf dem Gerät, Intervall bis 6 Monate einstellbar oder deaktivierbar & SkinVision, UMSkinCheck~\cite{app-umskincheck}, MoleScope & UC-10, UC-11 \\
\bottomrule
\end{tabularx}
\end{table}

\paragraph{Geführte Kameraaufnahme (F1).} Die App unterstützt beim Fotografieren, statt nur die Kamera zu öffnen. Ein Overlay im Kamerabild zeigt, wo die Hautstelle liegen und wie groß sie erscheinen soll. So findet der Nutzer den passenden Abstand und Bildausschnitt und löst das Foto anschließend selbst aus. Schärfe und Belichtung prüft die App danach (F2). So entstehen gleichmäßige Aufnahmen, die sich für die KI-Analyse und den späteren Vergleich im Verlauf eignen.

Der Hauptanwendungsfall UC-01 \emph{Hautstelle prüfen} verbindet die Features F1 bis F5 zu einem Ablauf.

\subsection{Nicht-MVP-Funktionalitäten}
\label{sec:nicht-mvp}

Die folgenden Funktionalitäten gehören zur langfristigen Vision und werden im Kurs nur umgesetzt, wenn Zeit bleibt:

\begin{table}[htbp]
\caption{Nicht-MVP-Funktionalitäten}
\label{tab:nicht-mvp}
\footnotesize
\renewcommand{\arraystretch}{1.15}
\begin{tabularx}{\textwidth}{@{}l >{\RaggedRight\arraybackslash}p{4.2cm} X@{}}
\toprule
\textbf{ID} & \textbf{Feature} & \textbf{Beschreibung} \\
\midrule
F9 & Arztsuche & Hautarzt in der Nähe über Google Maps finden \\
F10 & UV-Index-Warnung & Warnung bei hohem UV-Index \\
F11 & Symptom-Fragebogen & Symptome wie Juckreiz oder Blutung angeben \\
F12 & Dermatologen-Zweitprüfung & Ergebnis von einem Hautarzt prüfen lassen \\
F13 & PDF-Bericht & Ergebnis als PDF exportieren \\
F14 & Seitenvergleich & Läsionsfotos direkt nebeneinander vergleichen \\
F15 & Risikoprofil und Hauttyp & Risikoprofil und Hauttyp dauerhaft pflegen \\
F16 & Teledermatologie & Fall an einen Facharzt senden \\
F17 & Konto und Einwilligung & Registrieren und Einwilligungen verwalten \\
\bottomrule
\end{tabularx}
\end{table}

\subsection{Kriterien für die MVP-Auswahl}

Bei der Auswahl der MVP-Funktionen haben wir uns an vier Kriterien orientiert:

\begin{enumerate}
\item \textbf{Kernnutzen:} Trägt die Funktion direkt dazu bei, verdächtige Hautveränderungen zu erkennen?
\item \textbf{Sicherheit:} Verhindert die Funktion Fehleinschätzungen, etwa durch schlechte Bildqualität oder ungeeignete Fälle?
\item \textbf{Machbarkeit:} Lässt sich die Funktion mit unserem Wissensstand im Kurszeitraum umsetzen?
\item \textbf{Abhängigkeiten:} Bauen andere Funktionen darauf auf?
\end{enumerate}

\subsection{Machbarkeit im Kurszeitraum}
\label{sec:machbarkeit}

Bis zur Projektabgabe am 05.02.2027 bleiben ab Mitte November knapp 12 Wochen. In diesem Zeitraum ist der MVP-Umfang mit 4 Personen realistisch, wenn wir für die KI-Analyse auf ein vortrainiertes Modell zurückgreifen und es mit einem öffentlichen Datensatz wie HAM10000 oder dem ISIC-Archiv nachtrainieren. Das größte Risiko liegt in der Genauigkeit des Modells bei Fotos, die mit dem Smartphone aufgenommen wurden. Die 2D-Bodymap ist der aufwendigste Teil. Falls die Zeit knapp wird, ersetzen wir sie durch eine einfache Auswahl der Körperregion aus einer Liste.
